\scanpage{63}
% block: numerical data continued
and the class of $\operatorname{End}(\deleted{\unsure{S|B}}\handwritten{M})$ in $\operatorname{Br}(Z_{\handwritten{\sigma}})$ (the pair $(Z_{\handwritten{\sigma}},\handwritten{\xi_\sigma})$ being defined up to isomorphism). Without assuming $\handwritten{\mathcal{M}}$ semi\handwritten{\unsure{(}}-si\deleted{p}mple, one can\inserted{\handwritten{still\begin{tikzpicture}[x=1cm,y=1cm,baseline=-.1cm]\draw[thin] (0,.15).. controls (.1,-.1) and (-.2,-.1)..(-1,-.1)--(-2.1,-.1).. controls (-2.3,-.1) and (-2.35,-.2)..(-2.35,-.5);\end{tikzpicture}}} say that the largest semi-simple subcategory(gener\deleted{g}ated by the semi-simple objects) is known up to equiva\unsure{\deleted{e}l|\deleted{l}e}ence from the preceding knowledge, which therefore implies knowledge of $K(\handwritten{\mathcal{M}})=\handwritten{\mathbb{Z}^{(\Sigma)}}$,\inserted{\handwritten{its effective part,\begin{tikzpicture}[x=1cm,y=1cm,baseline=-.1cm]\draw[thin] (0,.15).. controls (.1,-.1) and (-.2,-.1)..(-1,-.1)--(-2.1,-.1).. controls (-2.3,-.1) and (-2.35,-.2)..(-2.35,-.5);\end{tikzpicture}}} and \unsure{\deleted{d}its|\deleted{s}da} variance for \unsure{\deleted{d}the|\deleted{l}} changes of\deleted{s} base by $k'$ varying over $k$. More precisely,$\handwritten{\Sigma'=\Sigma(\mathcal{M}_{k'})}$ as a set over $\handwritten{\Sigma}$ corresponds to the set \deleted{of fields} whose fibre \deleted{is} at $\handwritten{\sigma}$ is the set of component fields of $Z_{\handwritten{\sigma}}\otimes_k k'$, say $Z'^{i}_{\handwritten{\sigma}}\deleted{\unsure{i|\times}}$, the map $K(\handwritten{\mathcal{M}})=\underline{Z}^{(\handwritten{\Sigma})}\handwritten{\longrightarrow}K(\handwritten{\mathcal{M}_{k'}})=\underline{Z}^{(\handwritten{\Sigma'})}$ is given by
\[
\handwritten{\sigma\longmapsto}\ n_i\ (\handwritten{\sigma},i)\ ,
\]
where $n_i$ is given by the following rule: if $d$ is \deleted{the degree of} the order of $\handwritten{\xi_\sigma\in}\operatorname{Br}(Z_{\handwritten{\sigma}})$, that of its image in $\operatorname{Br}(Z'^{i}_{\handwritten{\sigma}})$ is a \deleted{idivs} divisor $d'_i$ of $d$, and we have $n_i=d/d_i$. (To be generalized to the case where the $Z_{\handwritten{\sigma}}$ are not necessarily separable). If $k'$ is an \deleted{\unsure{clo}} algebraically closed extension of $k$, the set $\handwritten{\Sigma'}$ is identified with the disjoint union of the \uline{fibers} at $k'$ of the $\operatorname{Spec}(Z_{\handwritten{\sigma}})$ ....

% block: b numerical characters
b) Suppose now $\handwritten{\mathcal{M}}$ \deleted{semi-Tannakian} equipped with a tensor product with associativity and commutativity constraints, with 1.

We designate by the name ``numerical characters\deleted{''} of $\handwritten{\Sigma}$ '' the following data:
\margin{\handwritten{make this notion precise, by making a ``numerical character'' into a structure capable of functorial variation....}}

% block: numerical characters item 1
\unsure{a|1}) the set (up to bijection)\inserted{\handwritten{$\deleted{\unsure{\xi|\Sigma}}\Sigma$\begin{tikzpicture}[x=1cm,y=1cm,baseline=-.1cm]\draw[thin] (-.1,-.25).. controls (0,0) and (.15,.1)..(.3,.4);\end{tikzpicture}}} of \deleted{com} isomorphism classes of simple objects of $\handwritten{\mathcal{M}}$.

% block: numerical characters item 2
2) The datum of the multiplication of $K(\handwritten{\mathcal{M}})={\unsure{\mathbb{Z}|B|Z}}^{(\handwritten{\Sigma})}$; which amounts to the datum for every $\handwritten{\sigma,\tau\in\Sigma}$ of the product
\[
\handwritten{\sigma\tau}=\handwritten{\sum_\rho}c_{\handwritten{\sigma\tau}}^{\handwritten{\rho}}\handwritten{\rho}\qquad\handwritten{(c_{\sigma,\tau}^{\rho}\in\mathbb{N})},
\]
\handwritten{i.e. of the system of natural numbers $\sigma_{\sigma,\tau}^{\rho}$.}
% archive sheet number: 62
